% from ega4-9-fr.tex
\oldpage[IV-3]{79}

élément de $\operatorname{Ass}(M_{(K)})$, et soit $\mathfrak p$ l'idéal premier
de $B$ image réciproque de $\mathfrak q$ par l'application canonique
$B\to B_{(K)}$. On sait (5.11.1.1) qu'il existe une filtration finie
$(N_h)_{0\leqslant h\leqslant m}$ de $M_{(K)}$ telle que $N_0=M_{(K)}$,
$N_m=0$, et que $N_h/N_{h+1}$ soit isomorphe à un sous-$B_{(K)}$-module
$\neq0$ de $B_{(K)}/\mathfrak q$. Soit $M_h$ l'image réciproque de $N_h$ par
l'application canonique $M\to M_{(K)}$. Il résulte de (9.4.4) que pour $s$
assez voisin de $\eta$, $(M_h)_{(\boldsymbol{k}(s))}$ s'identifie à un
sous-$B_{(\boldsymbol{k}(s))}$-module de $M_{(\boldsymbol{k}(s))}$, et le
quotient
$(M_h)_{(\boldsymbol{k}(s))}/(M_{h+1})_{(\boldsymbol{k}(s))}$ à un
sous-$B_{(\boldsymbol{k}(s))}$-module $\neq0$ de
$B_{(\boldsymbol{k}(s))}/\mathfrak p_{(\boldsymbol{k}(s))}
=(B/\mathfrak p)_{(\boldsymbol{k}(s))}$. Compte tenu de (3.1.7), on voit qu'on
est ramené à prouver, avec les mêmes notations, que si $B$ est \emph{intègre},
alors $B_{(\boldsymbol{k}(s))}$ n'a pas d'idéaux premiers associés immergés
pour $s$ voisin de $\eta$. Or, en remplaçant au besoin $A$ par $A_g$ et $B$ par
$B_g$ (où $g$ est un élément $\neq0$ de $A$), on peut supposer que $B$ contient
une $A$-algèbre de polynômes $C=A[\mathrm T_1,\ldots,\mathrm T_n]$ telle que
$B$ soit un $C$-module de type fini (Bourbaki, \emph{Alg. comm.}, chap. V,
\S~3, n$^\circ$~1, cor. 1 du th. 1). Comme $B_{(K)}$ est un $C_{(K)}$-module
sans torsion, on peut appliquer de nouveau le raisonnement fait plus haut en
remplaçant $B$, $M$ et $\mathfrak q$ par $C$, $B$ et $(0)$ respectivement, et
il suffit donc de voir que pour $s$ voisin de $\eta$, $C_{(\boldsymbol{k}(s))}$
n'a pas d'idéaux premiers associés immergés. Mais cela est évident puisque
$C_{(\boldsymbol{k}(s))}=\boldsymbol{k}(s)[\mathrm T_1,\ldots,\mathrm T_n]$
est un anneau intègre. C.Q.F.D.

\begin{theorem}[9.7.7]
\label{IV.9.7.7-fr}
Soit $f:X\to S$ un morphisme de présentation finie, et soit $E$ l'ensemble des
$s\in S$ pour lesquels $X_s$ a l'une des propriétés suivantes~: être~:
\begin{enumerate}[label=(\roman*)]
  \item géométriquement irréductible~;
  \item géométriquement connexe~;
  \item géométriquement réduit~;
  \item géométriquement intègre.
\end{enumerate}
Alors $E$ est localement constructible dans $S$.
\end{theorem}

Pour voir que les propriétés envisagées sont constructibles, remarquons d'abord
qu'elles vérifient trivialement la condition (9.2.1, (i)), en vertu de
(4.5.6, (i)) et (4.6.5, (i)). Il reste donc à vérifier (9.2.1, (ii)), donc on
peut supposer $S=\operatorname{Spec}(A)$ affine noethérien et intègre de point
générique $\eta$. En vertu de (4.6.8), il existe une extension finie $K'$ de
$K=\boldsymbol{k}(\eta)$ telle que $(X_\eta)_{(K')}$ soit tel que ses
composantes irréductibles (resp. connexes) soient géométriquement irréductibles
(resp. géométriquement connexes) et que $((X_\eta)_{(K')})_{\mathrm{red}}$ soit
géométriquement réduit. Comme il existe une base de $K'$ sur $K$ formée
d'éléments entiers sur $A$, l'anneau intègre $A'$ engendré par ces éléments est
fini sur $A$ et a $K'$ pour corps des fractions. Posons
$S'=\operatorname{Spec}(A')$~; le morphisme $g:S'\to S$ est fini et dominant,
donc surjectif (\textbf{II}, 6.1.10). Posons $X'=X_{(S')}$, de sorte que si
$\eta'$ est le point générique de $S'$, on a
$X'_{\eta'}=(X_\eta)_{(K')}$~; l'ensemble $E'$ des $s'\in S'$ tels que
$X'_{s'}$ ait l'une des propriétés (i), (ii), (iii) ou (iv) est égal à
$g^{-1}(E)$ (9.2.2, (iv)) ($E$ correspondant bien entendu à la même propriété)~;
$g$ étant surjectif, on a $E=g(E')$ et $S-E=g(S'-E')$~; en outre, $g$ est fermé
et $g^{-1}(\eta)=\{\eta'\}$ puisque $A'$ est intègre et fini sur $A$ (Bourbaki,
\emph{Alg. comm.}, chap. V, \S~2, n$^\circ$~1, cor. 1 de la prop. 1), donc
l'image par $g$ de tout voisinage de $\eta'$ est un voisinage de $\eta$. Le
théorème sera donc démontré si l'on prouve que $E'$ ou $S'-E'$ est un voisinage
de $\eta'$. Autrement dit, on peut désormais supposer que les composantes
irréductibles (resp. connexes) de $X_\eta$ sont géométriquement

